% Strana 1 titulní strany
% Napodobuje odstavce z DOCX: písmo Calibri Light, jednoduché řádkování
% (výška řádku ≈ 1,2264 × velikost písma — změřeno z exportu Wordu),
% 6 pt za odstavcem (styl Normální). Prázdný 14pt odstavec = 17,17 + 6 = 23,17 pt.
% Word klade účaří do vzdálenosti (spodek písma předchozího řádku + horní část
% písma dalšího): horní část = 0,952 × velikost, spodek = 0,269 × velikost.
% LaTeX bere jen \baselineskip dalšího řádku, proto korekce při změně velikosti:
%   20pt → 14pt: 5,38 + 13,33 − 17,09 = +1,62 pt
%   14pt → 20pt: 3,77 + 19,04 − 24,41 = −1,60 pt
\begin{titlepage}
    \thispagestyle{empty}
    \setlength{\parskip}{0pt}
    \setlength{\parindent}{0pt}
    % výšky řádků níže jsou přesné hodnoty z Wordu — bez globálního \linespread{1.15}
    \linespread{1}

    % Název školy: 20pt tučně, verzálky, centrováno.
    % Podpěra 42,13 pt = 1 prázdný odstavec nad ním + horní část 20pt písma;
    % účaří tak leží přesně jako ve Wordu, nezávisle na diakritice.
    % (\vspace* na začátku stránky by přidal ještě nežádoucí mezeru.)
    {\centering\headingfont\fontsize{20}{24.41}\selectfont\bfseries
     \rule{0pt}{42.63pt}\MakeUppercase{\varUniversity}\par}
    \vspace{6pt}

    % Studijní program: 14pt, řádkování 1,5 (25,76 pt), bez mezery za odstavcem.
    % Word dává přídavek řádkování 1,5 (8,59 pt) pod řádek, LaTeX nad něj → posun.
    \vspace{\dimexpr -8.59pt+1.62pt\relax}% + přechod 20pt → 14pt
    {\centering\headingfont\fontsize{14}{25.76}\selectfont\varProgram\par}
    \vspace{8.59pt}

    % 9 prázdných odstavců (9 × 23,17 pt) + přechod 14pt → 20pt
    \vspace{\dimexpr 208.53pt-1.60pt\relax}

    % Název práce: 20pt tučně, verzálky, centrováno (delší název se zalomí jako ve Wordu)
    {\centering\headingfont\fontsize{20}{24.41}\selectfont\bfseries
     \MakeUppercase{\varThesisName}\par}
    \vspace{\dimexpr 6pt+1.62pt\relax}% + přechod 20pt → 14pt

    % Typ práce: 14pt
    {\centering\headingfont\fontsize{14}{17.09}\selectfont\varThesisType\par}
    \vspace{6pt}

    % 10 prázdných odstavců (10 × 23,17 pt); u delšího názvu práce se mezera
    % zmenší, aby autor, vedoucí i místo zůstali na titulní straně
    \vspace{231.7pt minus 150pt}

    % Autor a vedoucí: 14pt, řádkování 1,5 (posun 8,59 pt jako u programu), zarovnáno vlevo,
    % tabulátor na 1985 twips = 99,25 pt od levého okraje textu
    \vspace{-8.59pt}
    {\headingfont\fontsize{14}{25.76}\selectfont
     \makebox[99.25pt][l]{Autor práce:}\varAuthor\par
     \makebox[99.25pt][l]{Vedoucí práce:}\varTutor\par}
    \vspace{8.59pt}

    % 2 prázdné odstavce (2 × 23,17 pt)
    \vspace{46.34pt}

    % Město a rok: 14pt, centrováno — hned pod vedoucím, ne u spodního okraje
    {\centering\headingfont\fontsize{14}{17.09}\selectfont\varCity\ \varYear\par}

\end{titlepage}
